%%% Cover Letter — Information Technology for Development (Taylor & Francis)
%%% Author: Tanzim Islam Khan

\documentclass[12pt,a4paper]{letter}

\usepackage[a4paper, margin=2.5cm]{geometry}
\usepackage[T1]{fontenc}
\usepackage[utf8]{inputenc}
\usepackage{mathptmx}
\usepackage{hyperref}
\usepackage{setspace}

\singlespacing

\signature{Tanzim Islam Khan}
\address{Tanzim Islam Khan\\
         dihan2468@gmail.com\\
         \today}

\begin{document}

\begin{letter}{The Editor-in-Chief\\
               \textit{Information Technology for Development}\\
               Taylor \& Francis}

\opening{Dear Editor-in-Chief,}

I am pleased to submit for your consideration the manuscript titled:

\begin{center}
\textbf{Generative AI Adoption in the Bangladeshi Software Engineering
Workforce: A Simulation-Based Analysis of Diffusion, Productivity,
and Skills Gaps}
\end{center}

\textbf{What the paper does.}
This paper presents the first integrated quantitative account of
generative AI (GenAI) adoption in the Bangladeshi software engineering
workforce between 2022 and 2026. The analysis combines a Bass diffusion
model calibrated to the adoption window, a reproducible synthetic
workforce simulation ($N = 1{,}500$) anchored to BASIS, BBS, and World
Bank population statistics, and a battery of statistical analyses
including heterogeneous treatment effect estimation by English
proficiency stratum and a skills gap index mapped against university
tier. Four policy scenarios are projected to 2031.

\textbf{Why it fits \textit{Information Technology for Development}.}
The journal's stated scope covers theoretical, empirical, and critical
research on the role of information technology in social, economic,
and human development outcomes in developing countries. This manuscript
addresses that scope directly on four dimensions:
\begin{itemize}
  \item \textit{Country focus:} the analysis is specific to Bangladesh,
    a large lower-middle-income country in which the IT sector is a
    primary vehicle for export-led growth and employment.
  \item \textit{Development relevance:} the findings are framed
    explicitly around SDG 8 (Decent Work and Economic Growth) and
    SDG 9 (Industry, Innovation and Infrastructure), both of which are
    central to Bangladesh's Vision 2041 development agenda.
  \item \textit{LMIC-specific insight:} the paper identifies English
    language proficiency as a robust structural moderator of GenAI
    productivity effects --- a finding specific to LMIC contexts that
    the dominant North American and European literature cannot supply.
  \item \textit{Policy utility:} the projected workforce scenarios and
    the skills gap index provide a quantitative basis for national
    digital literacy and university curriculum policy, directly usable
    by a2i, the ICT Division, and the University Grants Commission
    of Bangladesh.
\end{itemize}

\textbf{Originality.}
To the best of my knowledge, no prior study integrates a diffusion
model, a calibrated workforce simulation, and a policy projection
specific to Bangladesh. The paper fills this gap and provides a
replicable analytical framework that can be adapted to other
comparable LMIC contexts (Vietnam, Philippines, Nigeria, Egypt).

\textbf{Reproducibility.}
All source code (Python), 40 automated unit tests, 12 publication-quality
figures, and 11 data tables are released under open licences at:
\begin{center}
\url{https://github.com/tanzimkhan24/genai-se-bangladesh}
\end{center}

\textbf{Submission declaration.}
This manuscript has not been published elsewhere and is not currently
under review by any other journal or conference. All work is original.
No external funding was received for this research. The author declares
no conflicts of interest. This study uses simulated data calibrated to
publicly reported aggregate statistics; no human subjects research was
conducted and no ethics approval was required.

\textbf{Suggested reviewers.}
I suggest the following scholars as potential reviewers, given their
expertise in ICT4D, technology adoption in LMICs, and the Bangladeshi
IT sector:
\begin{enumerate}
  \item A researcher specialising in Bass diffusion modelling in
    emerging-economy ICT contexts.
  \item A researcher specialising in GenAI and software engineering
    productivity in South or Southeast Asia.
  \item A researcher specialising in computing education and skills
    gaps in LMIC university systems.
\end{enumerate}

I am happy to provide specific names on request, or to defer entirely
to the Editor's discretion.

I look forward to receiving the reviewers' comments and to the
possibility of publishing this work in \textit{Information Technology
for Development}.

\closing{Yours sincerely,}

\end{letter}
\end{document}
